\documentclass[../../main.tex]{subfiles}

\begin{document}

\section{Cantor 集和 Cantor-Lebesgue 函数}\label{section:ch02-Cantor}

我们已证明可数集测度为零以及 Borel 集是 Lebesgue 可测的. 这两个断言引发了以下两个问题.

\textbf{问题 1} 若一个集合测度为零, 它是否也可数?

\textbf{问题 2} 若一个集合可测, 它也是 Borel 集吗?

这两个问题的回答都是否定的. 本节我们构造称为 Cantor 集的集合与称为 Cantor-Lebesgue 函数的函数. 通过研究它们, 我们回答以上两个问题, 而后者提供了关于函数更好的性质的其他问题的回答.

考虑闭有界区间 $I=[0,1]$. 构造 Cantor 集的第一步是将 $I$ 分为长度等于 $1/3$ 的三个区间且去掉中间区间的内部, 即我们从区间 $[0,1]$ 去掉区间 $(1/3,2/3)$ 以得到闭集 $C_{1}$, 它是两个不交闭区间的并, 每个区间长度为 $1/3$:
\begin{align*}
C_{1}=[0,1/3]\cup[2/3,1].
\end{align*}
现在在 $C_{1}$ 的两个区间重复 "去掉居中的三分之一开集" 以得到闭集 $C_{2}$, 它是 $2^{2}$ 个闭区间的并, 每个区间的长度为 $1/3^{2}$:
\begin{align*}
C_{2}=[0,1/9]\cup[2/9,1/3]\cup[2/3,7/9]\cup[8/9,1].
\end{align*}
现在在 $C_{2}$ 的四个区间重复 "去掉居中的三分之一开集" 以得到闭集 $C_{3}$, 它是 $2^{3}$ 个闭区间的并, 每个区间的长度为 $1/3^{3}$.

继续这个去除操作可数多次以得到可数集族 $\{C_{k}\}_{k=1}^{\infty}$. 我们定义 Cantor 集 $\mathbf{C}$ 为
\begin{align*}
\mathbf{C}=\bigcap_{k=1}^{\infty}C_{k}.
\end{align*}

集族 $\{C_{k}\}_{k=1}^{\infty}$ 具有以下两个性质:
\begin{enumerate}[(i)]
\item $\{C_{k}\}_{k=1}^{\infty}$ 是下降的闭集序列;
\item 对每个 $k$, $C_{k}$ 是 $2^{k}$ 个不交闭区间的并集, 每个区间长度为 $1/3^{k}$.
\end{enumerate}

\begin{proposition}\label{proposition:ch02-Cantor零测}
Cantor 集 $\mathbf{C}$ 是闭的测度为零的不可数集.
\end{proposition}

\begin{proof}
任何闭集族的交是闭的. 因此 $\mathbf{C}$ 是闭的. 每个闭集是可测的. 因此每个 $C_{k}$ 和 $\mathbf{C}$ 自身是可测的.

现在每个 $C_{k}$ 是 $2^{k}$ 个不交区间的并, 每个区间长度为 $1/3^{k}$, 因此根据 Lebesgue 测度的有限可加性,
\begin{align*}
m(C_{k})=(2/3)^{k}.
\end{align*}
根据测度的单调性, 由于对所有 $k$, $m(\mathbf{C})\leqslant m(C_{k})=(2/3)^{k}$, $m(\mathbf{C})=0$. 剩下要证明 $\mathbf{C}$ 是不可数的. 为此我们用反证法讨论. 假定 $\mathbf{C}$ 是可数的. 令 $\{c_{k}\}_{k=1}^{\infty}$ 为 $\mathbf{C}$ 的一个列举. 其并是 $C_{1}$ 的两个不交 Cantor 区间之一不包含点 $c_{1}$, 记之为 $F_{1}$. 其并是 $F_{1}$ 的 $C_{2}$ 中两个不交 Cantor 区间之一不包含点 $c_{2}$, 记之为 $F_{2}$. 以这种方式继续, 我们构造了一个可数集族 $\{F_{k}\}_{k=1}^{\infty}$, 其中, 对每个 $k$, 它具有以下三个性质: (i) $F_{k}$ 是闭的且 $F_{k+1}\subseteq F_{k}$; (ii) $F_{k}\subseteq C_{k}$; (iii) $c_{k}\notin F_{k}$. 从 (i) 和集套定理我们得出交集 $\bigcap_{k=1}^{\infty}F_{k}$ 是非空的. 令点 $x$ 属于这个交. 根据性质 (ii),
\begin{align*}
\bigcap_{k=1}^{\infty}F_{k}\subseteq\bigcap_{k=1}^{\infty}C_{k}=\mathbf{C}.
\end{align*}
因此点 $x$ 属于 $\mathbf{C}$. 然而, $\{c_{k}\}_{k=1}^{\infty}$ 是 $\mathbf{C}$ 的列举, 使得对某个指标 $n$, $x=c_{n}$. 因此 $c_{n}=x\in\bigcap_{k=1}^{\infty}F_{k}\subseteq F_{n}$. 这与性质 (iii) 矛盾. 因此 $\mathbf{C}$ 必须是不可数的.
\end{proof}

定义在实数集上的实值函数 $f$ 称为是递增的, 若每当 $u\leqslant v$, $f(u)\leqslant f(v)$; 而称为是严格递增的, 若每当 $u<v$, $f(u)<f(v)$.

我们现在定义 Cantor-Lebesgue 函数, 即定义在 $[0,1]$ 上的连续的递增函数 $\varphi$, 它具有非凡的性质, 尽管 $\varphi(1)>\varphi(0)$, 它的导数存在且在测度为 $1$ 的集合上等于零. 对每个 $k$, 令 $\mathcal{O}_{k}$ 为在 Cantor 消去过程的前 $k$ 步中被去掉的 $2^{k}-1$ 个区间的并集. 因此 $C_{k}=[0,1]\sim\mathcal{O}_{k}$. 定义 $\mathcal{O}=\bigcup_{k=1}^{\infty}\mathcal{O}_{k}$. 则根据 De Morgan 等式, $\mathbf{C}=[0,1]\sim\mathcal{O}$. 我们首先在 $\mathcal{O}$ 上定义 $\varphi$, 接着在 $\mathbf{C}$ 上定义它.

固定自然数 $k$. 在 $\mathcal{O}_{k}$ 上定义 $\varphi$ 为 $\mathcal{O}_{k}$ 上的递增函数, 它在 $2^{k}-1$ 个开区间上是常数且取 $2^{k}-1$ 个值 $\{1/2^{k},2/2^{k},3/2^{k},\cdots,[2^{k}-1]/2^{k}\}$. 因此, 在消去过程的第一步被去掉的单个区间上, $\varphi$ 的定义是若 $x\in(1/3,2/3)$, $\varphi(x)=1/2$. 在前两步去掉的三个区间上, $\varphi$ 取 $1/4,2/4,3/4$.

通过在 $\mathbf{C}$ 上定义我们将 $\varphi$ 延拓到整个 $[0,1]$ 区间:
\begin{align*}
\varphi(0)=0,\qquad\varphi(x)=\sup\{\varphi(t)\mid t\in\mathcal{O}\cap[0,x)\},\quad\text{若 }x\in\mathbf{C}\sim\{0\}.
\end{align*}

\begin{proposition}\label{proposition:ch02-Cantor连续}
Cantor-Lebesgue 函数 $\varphi$ 是递增的连续函数, 它将 $[0,1]$ 映上 $[0,1]$. 它的导数在开集 $\mathcal{O}$ 上存在, 即 Cantor 集在 $[0,1]$ 上的补集存在,
\begin{align*}
\text{在 }\mathcal{O}\text{ 上 }\varphi'=0\quad\text{而 }m(\mathcal{O})=1.
\end{align*}
\end{proposition}

\begin{proof}
由于 $\varphi$ 在 $\mathcal{O}$ 上是递增的, 它在 $[0,1]$ 上的延拓也是递增的. 关于连续性, $\varphi$ 当然在 $\mathcal{O}$ 的每个点是连续的, 这是由于每个这样的点属于一个开区间, 在这个开区间上 $\varphi$ 是常数. 现在考虑点 $x_{0}\in\mathbf{C}$ 满足 $x_{0}\neq 0,1$. 由于点 $x_{0}$ 属于 $\mathbf{C}$, 它不是消去过程前 $k$ 步去掉的 $2^{k}-1$ 个区间的并, 其并我们记为 $\mathcal{O}_{k}$. 因此, 若 $k$ 充分大, $x_{0}$ 落在 $\mathcal{O}_{k}$ 中两个相连的区间之间, 选取其中小的为 $a_{k}$, 大的为 $b_{k}$. 根据函数 $\varphi$ 的定义, 穿过 $\mathcal{O}_{k}$ 中两个相连的区间增加 $1/2^{k}$. 因此
\begin{align*}
a_{k}<x_{0}<b_{k}\quad\text{且}\quad\varphi(b_{k})-\varphi(a_{k})=1/2^{k}.
\end{align*}
由于 $k$ 可以任意大, 函数 $\varphi$ 在 $x_{0}$ 没有跳跃不连续性. 对于一个递增的函数, 跳跃不连续性是仅有的不连续类型. 因此 $\varphi$ 在 $x_{0}$ 是连续的. 若 $x_{0}$ 是 $[0,1]$ 的端点, 用类似的方法可证明在 $x_{0}$ 的连续性.

由于在每个去除过程的任意阶段去掉的区间上 $\varphi$ 是常数, 它的导数存在且在 $\mathcal{O}$ 的每一点等于 $0$. 由于 $\mathbf{C}$ 的测度为零, 它在 $[0,1]$ 上的补 $\mathcal{O}$ 的测度为 $1$. 最后, 由于 $\varphi(0)=0$, $\varphi(1)=1$, 而 $\varphi$ 是递增与连续的, 我们从介值定理推出 $\varphi$ 将 $[0,1]$ 映上 $[0,1]$.
\end{proof}

\begin{proposition}\label{proposition:ch02-Cantor映射}
令 $\varphi$ 为 Cantor-Lebesgue 函数且定义在 $[0,1]$ 上的函数 $\psi$ 为
\begin{align*}
\text{对所有 }x\in[0,1],\quad\psi(x)=\varphi(x)+x.
\end{align*}
则 $\psi$ 是严格递增的连续函数, 它将 $[0,1]$ 映上 $[0,2]$.
\begin{enumerate}[(i)]
\item 将 Cantor 集 $\mathbf{C}$ 映上正测度的可测集;
\item 将一个可测集即 Cantor 集的一个子集映上一个非可测集.
\end{enumerate}
\end{proposition}

\begin{proof}
由于函数 $\psi$ 是两个连续函数的和, 它是连续的; 由于它是递增与严格递增函数的和, 它是严格递增的. 此外, 由于 $\psi(0)=0$ 且 $\psi(1)=2$, $\psi([0,1])=[0,2]$. 对于 $\mathcal{O}=[0,1]\sim\mathbf{C}$, 我们有不交分解
\begin{align*}
[0,1]=\mathbf{C}\cup\mathcal{O},
\end{align*}
其中 $\psi$ 有不交分解
\begin{align}
[0,2]=\psi(\mathcal{O})\cup\psi(\mathbf{C}).\label{align:ch02-psi分解}
\end{align}

定义在区间上的严格递增函数有一个连续的逆. 因此 $\psi(\mathbf{C})$ 是闭的而 $\psi(\mathcal{O})$ 是开的, 于是它们都是可测的. 我们将证明 $m(\psi(\mathcal{O}))=1$, 且因此从\eqref{align:ch02-psi分解}推出 $m(\psi(\mathbf{C}))=1$, 从而证明 (i).

令 $\{I_{k}\}_{k=1}^{\infty}$ 为在 Cantor 去除过程去掉的区间族（以任何方式）的列举. 因此 $\mathcal{O}=\bigcup_{k=1}^{\infty}I_{k}$. 由于 $\varphi$ 在每个 $I_{k}$ 上是常数, $\psi$ 将 $I_{k}$ 映上它自身的具有相同长度的平移副本. 由于 $\psi$ 是一对一的, 族 $\{\psi(I_{k})\}_{k=1}^{\infty}$ 是不交的. 根据测度的可数可加性,
\begin{align*}
m(\psi(\mathcal{O}))=\sum_{k=1}^{\infty}\ell(\psi(I_{k}))=\sum_{k=1}^{\infty}\ell(I_{k})=m(\mathcal{O}).
\end{align*}
但 $m(\mathbf{C})=0$ 从而 $m(\mathcal{O})=1$. 因此 $m(\psi(\mathcal{O}))=1$, 且根据\eqref{align:ch02-psi分解}, $m(\psi(\mathbf{C}))=1$. 我们证明了 (i).

为了证明 (ii), 注意到 Vitali 定理告诉我们 $\psi(\mathbf{C})$ 包含不可测集 $W$. 集合 $\psi^{-1}(W)$ 是可测的且测度为零, 因为它是 Cantor 集的子集. 集合 $\psi^{-1}(W)$ 是 Cantor 集的可测子集, 它被 $\psi$ 映上一个不可测集.
\end{proof}

\begin{proposition}\label{proposition:ch02-Cantor非Borel}
存在 Cantor 集的一个子集是可测集但不是 Borel 集.
\end{proposition}

\begin{proof}
前一个命题描述的定义在 $[0,1]$ 上的严格递增函数 $\psi$ 将可测集 $A$ 映上一个不可测集. 定义在区间上的严格递增连续函数将 Borel 集映上 Borel 集（见习题 43）. 因此集合 $A$ 不是 Borel 集, 否则它在 $\psi$ 下的象将是 Borel 集, 因此是可测的.
\end{proof}

\begin{exercise}\label{exercise:ch02-34}
证明存在 $[0,1]$ 上的连续、严格递增函数将正测度集映到零测度集.
\end{exercise}

\begin{exercise}\label{exercise:ch02-35}
令 $f$ 为开区间 $I$ 上的递增函数. 对于 $x_{0}\in I$, 证明 $f$ 在 $x_{0}$ 连续当且仅当存在 $I$ 中的序列 $\{a_{n}\}$ 和 $\{b_{n}\}$ 使得对每个 $n$, $a_{n}<x_{0}<b_{n}$ 且 $\lim_{n\to\infty}[f(b_{n})-f(a_{n})]=0$.
\end{exercise}

\begin{exercise}\label{exercise:ch02-36}
令 $f$ 为定义在 $E$ 上的连续函数. 若 $A$ 是可测的, $f^{-1}(A)$ 总是可测的吗?
\end{exercise}

\begin{exercise}\label{exercise:ch02-37}
令函数 $f:[a,b]\to\mathbf{R}$ 为 Lipschitz 的, 即存在常数 $c\geqslant 0$ 使得对所有 $u,v\in[a,b]$, $|f(u)-f(v)|\leqslant c|u-v|$. 证明 $f$ 将测度为零的集映上测度为零的集. 证明 $f$ 将 $F_{\sigma}$ 集映上 $F_{\sigma}$ 集. 证明 $f$ 将可测集映上可测集.
\end{exercise}

\begin{exercise}\label{exercise:ch02-38}
令 $F$ 为按与 Cantor 集相同的方式构造的 $[0,1]$ 的子集, 除了在第 $n$ 步每个被去掉的区间的长度为 $a3^{-n}$, 其中 $0<a<1$. 证明 $F$ 是闭集, $[0,1]\sim F$ 在 $[0,1]$ 稠密, 且 $m(F)=1-a$. 这样的集合 $F$ 称为广义 Cantor 集.
\end{exercise}

\begin{exercise}\label{exercise:ch02-39}
证明存在实数的开集有一个正测度的边界（与直觉相反）.（提示: 考虑前一个习题中的广义 Cantor 集的补集.）
\end{exercise}

\begin{exercise}\label{exercise:ch02-40}
$\mathbf{R}$ 的子集 $A$ 称为在 $\mathbf{R}$ 无处稠密, 若每个开集 $\mathcal{O}$ 有一个与 $A$ 不交的非空开子集. 证明 Cantor 集在 $\mathbf{R}$ 无处稠密.
\end{exercise}

\begin{exercise}\label{exercise:ch02-41}
证明定义在区间上的严格递增函数有一个连续的逆.
\end{exercise}

\begin{exercise}\label{exercise:ch02-42}
令 $f$ 为连续函数而 $B$ 是 Borel 集. 证明 $f^{-1}(B)$ 是 Borel 集.（提示: 使得 $f^{-1}(E)$ 是 Borel 集的集合 $E$ 所成的族是包含开集的 $\sigma$ 代数.）
\end{exercise}

\begin{exercise}\label{exercise:ch02-43}
用前两个习题证明定义在区间上的严格递增连续函数将 Borel 集映到 Borel 集.
\end{exercise}

\end{document}